% Generated by revision/scripts/calibration_study.py
\begin{table}[!h]
\centering\small
\caption{PathXDRP uncertainty quality on all five splits, mean $\pm$
standard deviation over five seeds. ECE measures \emph{absolute}
calibration; the selective-RMSE gain and the AUROC of $\sigma$ against
the large-error half of the test set measure the \emph{ranking} quality
that selective prediction actually depends on. The random split
alone is the most favourable of the five, which is why all five
are shown.}
\label{tab:calibration_all_splits}
\begin{tabular*}{\columnwidth}{@{\extracolsep{\fill}}lcccc@{}}
\toprule
\textbf{Split} & \textbf{ECE} $\downarrow$ & \textbf{sel.\ RMSE@50\%} $\downarrow$ & \textbf{gain} $\uparrow$ & \textbf{AUROC}$(\sigma)$ $\uparrow$ \\
\midrule
random & $0.213{\pm}0.048$ & $0.771$ & $24.2\%$ & $0.648$ \\
cell-blind & $0.148{\pm}0.108$ & $1.094$ & $23.8\%$ & $0.637$ \\
drug-blind & $1.428{\pm}0.249$ & $2.238$ & $17.8\%$ & $0.542$ \\
scaffold-blind & $1.298{\pm}0.339$ & $2.426$ & $2.7\%$ & $0.526$ \\
tissue-blind & $0.361{\pm}0.095$ & $1.205$ & $21.0\%$ & $0.617$ \\
\bottomrule
\end{tabular*}

\smallskip
\footnotesize\emph{Note:} `gain' is the relative reduction in RMSE
obtained by discarding the least-confident half of the predictions.
AUROC$(\sigma)$ is the area under the ROC curve for using $\sigma$ to
identify the half of the test set with the larger absolute residual;
$0.5$ is chance.
\end{table}